\clearpage
\section*{How accurate is the permit construction data, and what is one row?}
\textit{September 23, 2026. Agent-drafted data review; the unit-of-observation choice is Jacob's.}

A stratified random sample of 225 cases from the permit-based construction data (dataset at commit
\texttt{9736e3c8}; \texttt{tasks/audits/construction\_hand\_checks}) was reviewed case by case against permit text
and Assessor histories by agent reviewers, with the main claims checked against the data. Weighted by stratum
population, 90.7 percent of density-eligible buildings were correct. Buildings measured on the permit's own parcels,
three quarters of the sample, were right in 24 of 25 cases, as were lot links; the weakest group was the
Assessor-only 2006--2007 buildings (17 of 25). Among withheld cases, 18 of 25 ``no permit found'' buildings were
genuine new construction, 14 of them with a permit that should have linked; 21 of 25 unit disagreements were
recoverable; and of permits that reached no building, 13 of 25 were never built, 5 were tax-exempt housing without
Assessor characteristics, and only 2 were missed links.

One mechanism explained the largest error in four strata: when a condominium declaration or subdivision retires a
parcel number, the building reappears on new parcels. The permit kept its parent lot's record while the finished
homes sat unlinked, or an Assessor-only parent card duplicated its children. On the Cleveland--Hobbie--Oak site,
eight permits for rowhome rows and two condominium buildings all listed one parent lot; the Assessor recorded a
39-unit card on it in 2014 and about 48 child parcels from 2016. The paper's data counts this site 40 times (the
parent card and 39 homes); ours had counted it zero times in the density sample. Jacob chose one row per building.
Each permit's parcels now include the successors of its listed parcels (each new parcel descends from the nearest
retired parcel on its block), a record reached by several permits belongs to the permit at its address, and
hand-named lots and homes join their permit's parcels. That site is now eight building rows; six match their
permits' dwelling counts exactly, and two (449 W Hobbie, 1001 N Cleveland) reach only one of their homes and are
withheld. Allowing several parents per new parcel recovered one more row but broke others (94 fewer eligible
buildings) and was not adopted. Two hand decisions that had assigned a whole divided site to one permit were retired.

The file now has 18,051 rows (18,875 before) and 14,053 density-eligible buildings (14,033): 11,776 permitted
(11,564) and 2,277 Assessor-only (2,469; the difference is mostly parent cards removed as duplicates). Counted as
the paper does, by completion year 2006--2022, there are 13,022 eligible buildings against the paper's roughly
13,600, part of which is the paper's double counting of divided sites. ``No permit found'' buildings fell from
1,776 to 1,077; flagged measured buildings rose from 842 to 959, mostly divided sites whose successors include a
neighbor built without its own permit. The validation sample has not yet been re-reviewed on the rebuilt data, and
first-year records that are partial (duplicate cards, a pre-declaration transitional record) remain the main
measurement error. No regressions have been run on these data.
